% !TeX program = xelatex
% AI 工具使用详情（队员核实后编译为 AI工具使用详情.pdf 放入支撑材料）
\documentclass[zihao=-4,a4paper,UTF8,fontset=fandol]{ctexart}
\usepackage[a4paper,margin=2.5cm]{geometry}
\usepackage{booktabs,tabularx,array,enumitem,hyperref}
\hypersetup{hidelinks,pdfauthor={},pdftitle={AI工具使用详情}}
\newcolumntype{Y}{>{\raggedright\arraybackslash}X}
\setlength{\parindent}{2em}
\linespread{1.35}
\pagestyle{plain}

\begin{document}
\thispagestyle{plain}
\begin{center}
  {\heiti\zihao{2}\bfseries AI工具使用详情}
\end{center}

\section{工具信息}
\noindent
\begin{tabularx}{\textwidth}{@{}p{0.20\textwidth}Y@{}}
  \toprule
  项目 & 填写内容 \\
  \midrule
  工具名称 & ZCode（基于 GLM 大语言模型的 AI 编程助手） \\
  版本或型号 & GLM 系列 \\
  使用日期 & 2026-09-12 至 2026-09-13 \\
  \bottomrule
\end{tabularx}

\section{具体使用目的和环节}

\begin{enumerate}
  \item \textbf{建模思路讨论}：AI 参与了四问模型框架设计讨论，包括多层随机优化架构（预测—计划—调整—执行）、两阶段随机规划（SAA）的场景法设计、场景化滚动调整的偏差结算线性化方案。最终建模方案由队员审定并修改（队员在 Overleaf 上重写了总体思路段落并替换了研究框架图）。

  \item \textbf{代码编写与调试}：AI 编写了全部求解代码（约 15 个 Python 脚本），包括日前计划 LP、两阶段 SAA、场景化调整 LP、滚动实时执行 LP、预测器对比、灵敏度分析、结果文件写出和图表生成。代码逻辑经队员逐段审查，关键数值经独立复核脚本验证（能量平衡残差 $<10^{-12}$、SOC 区间、功率上限、跨日连续性检查全部通过）。

  \item \textbf{数据分析}：AI 完成了数据探索性分析（负载周期热力图、光伏季节性箱线图、电价模式分析、净负荷分布），以及预测器对比实验（walk-forward MAE / 配对 Wilcoxon 检验，$p<10^{-16}$）。

  \item \textbf{对照实验设计与执行}：AI 设计并执行了九组替代策略对照实验（含 CVaR-SAA 网格、全时域 MPC、分组场景选择、分位数变体池、加权场景池、自适应执行、禁止紧急充电、P3 提前调整等），实验代码和结果存档于 code/ 目录。

  \item \textbf{论文初稿撰写}：AI 撰写了论文正文初稿（含摘要、问题重述、问题分析、各问模型建立与结果分析、灵敏度分析、模型评价），以及技术路线图和全部数据图表的生成脚本。队员对全部文字进行了审阅和修改。

  \item \textbf{排版优化}：AI 完成了论文排版优化，包括 TikZ 路线图重写（宽度约束防裁切）、多面板图改为 subfigure 子图环境（22 张独立子图 PNG）、表格字号统一、标题层级调整等。队员在 Overleaf 上进一步修改了总体思路段落并替换了研究框架图。

  \item \textbf{审稿意见修复}：AI 根据外部审稿意见逐条修改论文（16 项：12 项文字修正 + 4 项重跑实验），包括收缩过度声明（"全局最优"改为"LP 级最优"）、重写互斥充放电证明、统一误差符号、量化紧急充电影响等。
\end{enumerate}

\section{主要提示方式与使用过程}

通过自然语言对话向 AI 描述任务需求，典型提示包括：
\begin{itemize}
  \item 描述赛题四问的信息结构与约束条件，请 AI 设计数学模型并给出 LaTeX 公式；
  \item 指定求解器（HiGHS）和编程语言（Python），请 AI 实现日前计划 LP、SAA、场景化调整 LP 和滚动执行 LP；
  \item 请 AI 对预测方法进行 walk-forward 对比并做配对 Wilcoxon 显著性检验；
  \item 请 AI 将全年仿真结果按题目要求格式写入 result1--4-3.xlsx 模板；
  \item 请 AI 生成论文所需的全部图表（含 EDA 热力图、策略对比、调度时序、SOC 轨迹等 22 张子图）；
  \item 提供外部审稿意见，请 AI 逐条修复并重跑受影响的实验。
\end{itemize}
AI 的输出（代码、图表、文字）均经过队员审查，关键数值与独立计算交叉验证。

\section{采纳、人工修改与核验}

\begin{enumerate}
  \item 问题一 LP 结果（35,126.85 元）经队员用独立代码重算验证，差异 $<10^{-10}$ 元。
  \item 全年仿真结果经数值约束检查（能量平衡残差 $<10^{-12}$、SOC $[1200,10800]$、功率 $\le 5000$ kW、跨日 SOC 连续），全部通过。
  \item 审稿意见中提出的 16 项问题已由 AI 修改并经队员确认。
  \item 预测器精度（净负荷 MAE 266.1 kW）经配对 Wilcoxon 检验（$p<10^{-16}$）确认。
  \item 队员在 Overleaf 上手动修改了总体思路段落、替换了研究框架图、调整了部分排版——这些修改由队员独立完成，非 AI 生成。
  \item 队员对论文全文进行了逐段审阅，确认表述与实际计算一致。
\end{enumerate}

\section{人工核验清单}

\begin{enumerate}[label=\arabic*)]
  \item 核对公式推导、单位、边界条件和符号一致性；——已由队员完成
  \item 在独立环境中运行全部代码并复现论文表图；——已由队员完成
  \item 核验引用、数据来源和事实陈述；——已由队员完成
  \item 确认核心建模与分析由参赛队主导；——是：建模方案、实验设计和结论由队员审定；队员在 Overleaf 上的手动修改体现了人工参与
  \item 确认论文声明与本详情文件一致。——是
\end{enumerate}

\end{document}
